\section{Method}
\subsection{Design and epistemological position}
We propose individual semi-structured interviews with a contextualist, experiential orientation. Accounts are treated as situated interpretations of experience, shaped by circumstances, platform affordances and the interview interaction. They are neither transparent measurements of underlying mental states nor merely linguistic objects detached from experience. Reflexive thematic analysis is appropriate to the aim of developing patterns of shared meaning across cases \citep{braun2021,byrne2022}. AIM-ER is a sensitising framework; analysis will remain open to meanings not captured by its constructs.

\subsection{Eligibility and recruitment}
Eligible participants will be adults aged 18 or older who can undertake an interview in Polish, have personally used a generative music system to create or revise music within the previous three months, and can provide informed consent. No formal musical training, positive experience, mental-health diagnosis, paid subscription or public portfolio is required. Listening to AI-generated songs without any creative participation does not meet the criterion. A session is an attempt to generate or revise music using a generative tool, bounded by the participant's reported start and stopping point; a completed song, public sharing and emotional motivation are not required. A session may end without a usable output. Conventional audio editing alone does not meet the criterion. Recruitment will use neutral invitations through several online creator communities and general-interest channels where permission has been obtained. Invitations will explicitly welcome people who discontinued use within the three-month eligibility window or found it unhelpful. People whose last creative session predates the window are outside this protocol; their experiences require a separate sampling design.

A short screening form will collect age eligibility, preferred contact route, recency of creation, self-described musical experience and frequency of use. Contact details will be kept outside the analysis dataset. Purposeful selection will seek variation in experience, platform use and perceived outcomes, rather than quotas implying representativeness. We will avoid recruitment of students dependent on the interviewer's grading or assessment, and prevent administrators from learning who declined. No recruitment or pilot interviews will begin before the applicable institutional review and permissions.

\subsection{Sampling rationale}
The provisional planning range is 20--30 completed interviews of approximately 45--60 minutes. This is a pragmatic range, not a power calculation or a guarantee of saturation. A relatively focused aim, recent episode recall and an existing theoretical framework can improve informational richness; diversity of participants and uneven interview depth may require further exploration \citep{malterud2016}. In line with reflexive thematic analysis, we will not use a count of newly generated codes as a mechanical stopping criterion \citep{braunsaturation2021}. After approximately 15 and 20 completed interviews the team will reflect on whether contrasting accounts, context and depth are adequate for the research questions and whether additional recruitment is feasible. At each review consider the focus of the questions, relevance of the sample, use of theory, depth of dialogue and intended cross-case analysis as qualitative information-power considerations, not a score. A decision to stop or extend within the approved range will be documented with reasons, including constraints. Count pilot interviews within the range only when included in the main corpus; report pilot and main recruitment separately. Recruitment beyond that range requires appropriate protocol amendment.

\begin{table}[htbp]
\centering\small
\caption{Purposeful variation dimensions, not mutually exclusive sampling strata.}
\label{tab:sampling}
\begin{tabularx}{\linewidth}{@{}p{.30\linewidth}X@{}}
\toprule
Dimension & Intended contrast \\
\midrule
Musical experience & No formal training; self-taught; trained or professionally experienced. \\
Generator use & Recent beginners; regular creators; people who reduced or stopped use. \\
Reported experience & Beneficial, neutral, disappointing, ambivalent or unpleasant episodes. \\
Creative operations & Minimal prompting; extensive selection or revision; hybrid manual and AI workflows. \\
Social context & Private creation; sharing with selected others; public-facing creation. \\
\bottomrule
\end{tabularx}
\end{table}
These dimensions will inform recruitment reflection, not hypotheses about subgroups. Participants can decline screening questions not required for eligibility.

\subsection{Pilot and fieldwork decision rules}
The two or three pilot interviews will assess the feasibility of the interview procedure, not the adequacy of the eventual thematic analysis. The interviewer and a second team member will review whether eligibility and episode selection were understood, whether the core domains could be discussed without unnecessary sensitive disclosure, whether prompts elicited concrete examples without implying a desired emotional outcome, whether the interview fitted the planned duration, and whether recording, consent and withdrawal procedures worked as approved. The review will distinguish a misunderstanding of the question from a legitimate refusal to answer. Neither refusal to discuss a private event nor absence of emotional change constitutes a pilot failure.

Before main recruitment, record one of three decisions: retain the guide with reasons; revise and review the changed domains; or suspend collection until ethical, technical or methodological problems are resolved. Material changes to approved procedures require the relevant institutional review. A pilot interview may enter the analytic corpus only if consent covers that use and the final research question, eligibility, interview domains and ethical procedures remain compatible. Document the decision for each pilot without publishing participant-level details. Pilot approval is not evidence that subsequent themes are valid.

At the reviews after approximately 15 and 20 completed interviews, consider whether accounts are sufficiently contextualised to address the questions, whether particular experiences are absent because of recruitment or interviewing practices, and whether additional interviews could contribute materially to interpretation. Seek relevant contrasts through approved recruitment channels rather than imposing numerical quotas. If a substantial gap remains at the approved maximum, report the limitation or seek an amendment; do not claim saturation merely because recruitment ended.

\subsection{Interview procedure}
Interviews will take place through an institutionally approved online service. Before the interview, participants receive a plain-language study description, information about recording and confidentiality limits, and an explanation of the right to skip questions or stop. Separate permission will be sought for audio recording. A researcher will conduct a brief technical check and reconfirm consent. No new song generation or disclosure of a traumatic event is required.

The primary unit of elicitation is the most recent eligible creative episode that the participant is willing and able to discuss, subject to the selection procedure below. The unit of analysis is the interview corpus; episodes are contextual units within accounts, not independent observations. The interviewer asks the participant to describe what they were doing before the session, what they intended, the sequence of actions, how they evaluated candidate outputs and what happened after stopping. An optional contrasting episode will help probe differences in motivation, process and perceived outcome. The interview moves from open recall to concrete prompts and only then to more specific questions about experience. The guide will be piloted with two or three eligible volunteers after approval, with changes recorded following the development principles of \citet{kallio2016}. Pilot material will enter the main corpus only if this is consented to and methodological compatibility is documented.

Participants may voluntarily describe an artefact or show a non-sensitive excerpt without transferring a file. The primary evidence remains the interview. Showing is not consent to capture, receive or retain an artefact. Pause recording and transcription before any live demonstration; resume only after the artefact is no longer visible or audible and the participant reconfirms consent. No screenshots, links, downloads or copied lyrics are retained by default. Separate approved consent and rights assessment are required for any proposed transfer, recording or retention; refusal has no bearing on inclusion. An interviewer will distinguish a participant's own description from an interpretation prompted by a lyric or audio example. For each episode, record broad tool type/version if known, intended use, temporal order and uncertainties in recall. The optional contrast is not a matched experimental comparison.
\subsection{Additional qualitative research domains}
The episode-centred guide will include optional follow-up prompts about the transition from lived experience to a musical representation: identifying or borrowing emotional labels, use of metaphor, changes in the narrative during revision, recognition of one's intended meaning, and the difference between first-person immersion and an observer perspective. Ask whether participants experienced insight, tension release, a sense of closure, pride, aesthetic pleasure or absorption, and whether such experiences persisted, changed or disappeared after later listening. Invite accounts of unchanged or worsened affect and of attempts that did not produce a usable artefact. Terms such as catharsis and self-therapy will be explored as participants' own interpretations rather than suggested benefits or clinical categories.

Ask about private creation, selective sharing and public posting, anticipated judgement, perceived anonymity, belonging, acceptance, rejection and unwanted disclosure. Explore whether creative activity facilitated subsequent communication or displaced it, and how participants decided to continue, revise, stop or switch strategies. Contrast repetitive listening or generation that felt purposeful with episodes described as difficult to disengage from. Where voluntarily offered, a participant may describe or privately present an artefact under the existing consent procedure; this is a recall aid, not independent evidence of prior affect. Do not elicit traumatic detail, assume positive effects, or infer a diagnosis from a song. Use these domains as sensitising prompts after free recall rather than a fixed coding scheme. The resulting interpretations may inform P03 and item content in P05, but do not validate scale structure or causal mechanisms.

\subsection{Candidate probes for contrasting episodes}
To operationalise the expanded interview scope, the following neutral probes supplement rather than replace the approved episode guide: What did you want to achieve before you started? Which decisions did you make yourself, and which did you delegate? Did the words or music represent something you already understood, or did they help you notice something new? Did the experience look different after listening, and if so, how? What happened when you returned to the song later? If another person heard the result, how did you interpret their response? What made you decide to stop or continue? Ask for counterexamples and experiences without any change, rather than eliciting agreement with a putative therapeutic mechanism.

Contrast a private episode with a shared episode only when the participant has such experiences and wishes to discuss them. A report of relief accompanied by postponed action is not contradictory; explore how both meanings coexist in the account. The differentiation of adaptive and maladaptive musical engagement in \citet{carlson2015} and the regulatory flexibility account of \citet{bonanno2013} inform attention to such contrasts, without imposing diagnostic categories. The analysis will use reflexive thematic interpretation \citep{braun2021}; a thematic matrix may organise cases by context and contrasting meanings but is not a frequency estimate or a reliability test of a predetermined set of mechanisms. The expanded prompts remain candidate guide revisions pending pilot and ethics review.

\subsection{Episode selection and interview comparability}
First invite the latest eligible session within three months, including incomplete or abandoned attempts, without asking for details before the participant decides whether to discuss it. If the person declines or cannot recall it sufficiently, invite the next most recent eligible session they are willing and able to discuss. Do not request the reason; optionally record only participant-approved categories (privacy, recall, other, not stated) and whether the latest session was discussed. The participant may decline every episode or stop without consequence. This procedure reduces selection of exemplary successes but does not eliminate comfort-related, recall or participation selection; the corpus cannot estimate the distribution of all sessions. Where feasible, ask for a contrasting episode selected by the participant, but do not require a second session. Record whether the account concerns a single session or repeated activity, approximate time elapsed, and which details the participant is uncertain about. Temporal order is reconstructed from recollection, not verified by timestamps. Figure~\ref{fig:episode-elicitation} separates the event sequence from the present interpretation and optional contrasting account. For each phase, distinguish what the participant remembers happening, what they recall thinking or feeling at the time, and what they now make of it. Ask whether later listening, sharing or another event intervened before a reported consequence; preserve uncertainty when the order cannot be reconstructed.

\begin{figure}[htbp]
\centering
\begin{tikzpicture}[
 phase/.style={draw,rounded corners=2pt,text width=2.65cm,minimum height=1.3cm,align=center,font=\small,fill=blue!5},
 note/.style={draw,rounded corners=2pt,text width=5.8cm,minimum height=1.25cm,align=center,font=\small,fill=black!3},
 arr/.style={-{Stealth[length=2mm]},thick},
 interpret/.style={-{Stealth[length=2mm]},dashed},
 node distance=3mm]
\node[phase] (a) {Before session\\Context and intention};
\node[phase,right=of a] (b) {Creation\\Actions and decisions};
\node[phase,right=of b] (c) {Stopping\\Output or no output};
\node[phase,right=of c] (d) {Later\\Listening, sharing, other events};
\draw[arr] (a)--(b); \draw[arr] (b)--(c); \draw[arr] (c)--(d);
\node[note,below=13mm of b,xshift=-1.5cm] (e) {Present interpretation\\Reported meaning, uncertainty, alternative explanations};
\node[note,right=5mm of e] (f) {Optional contrasting episode\\Similarities and differences in context, action and experience};
\draw[interpret] (a.south) -- (e.north west);
\draw[interpret] (c.south) -- (e.north east);
\draw[interpret] (d.south) -- (f.north);
\draw[interpret] (f.west)--(e.east);
\node[below=4mm of e,anchor=north west,font=\scriptsize,align=left,text width=12cm,xshift=-2cm] {Solid arrows: elicitation order. Dashed arrows: interpretive comparison.\\Neither arrow type identifies causation.};
\end{tikzpicture}

\caption{Prospective elicitation structure. The upper sequence is a recall scaffold, not an assumed causal model; steps may repeat, be absent or occur in another order. The lower boxes distinguish retrospective interpretation and optional contrast. No participant observations are represented.}
\label{fig:episode-elicitation}
\end{figure} The interviewer will separate the original intention from meanings assigned retrospectively, distinguish the session ending from later listening or sharing, and ask about other circumstances that may have contributed to changes. The absence of a reported change is a legitimate account.

A common core guide and documented optional probes permit comparison without requiring identical wording. The interviewer will not show the AIM-ER pathway diagram before open recall. After the pilot, retain the guide version, a brief rationale for revisions, and a record of departures from the common domains. The research team will assess feasibility in terms of interview depth, participant burden and capacity to conduct careful analysis, rather than setting a throughput target that compromises quality.
\subsection{Interview guide}
The table presents an English summary of the interview domains and illustrative prompts. The accompanying Polish guide specifies the wording for fieldwork. Follow-up questions will use the participant's language and ask for examples, counterexamples and clarification without suggesting benefit or harm.

\begin{table}[htbp]
\centering\small
\caption{Interview domains and neutral elicitation prompts.}
\label{tab:interview}
\begin{tabularx}{\linewidth}{@{}p{.23\linewidth}X@{}}
\toprule
Domain & Illustrative prompt \\
\midrule
Context and motive & Think of your latest eligible occasion, including an unfinished attempt. Would you be willing and able to discuss it? If not, we can discuss the next most recent occasion or skip this part. What led to that occasion? Did your intention change as the session progressed? \\
Translation & How did you decide what to put into the prompt, lyrics or musical settings? What did you leave out? \\
Iteration & What happened after the initial attempt, whether or not it produced an output? What led to changes, acceptance or stopping, including stopping without a usable output? \\
Contribution & Which aspects did you regard as your own work, and which as the system's? How did you experience that division and your ability to influence the process? If an artefact exists, to what extent did it feel like yours? What contributed to that perception? \\
Personal meaning & Did the result relate to your personal experience, if at all? How did it relate or fail to relate? Did it match your intentions? What indicated that? If no usable result exists, what significance, if any, did the attempt itself have? \\
Consequences & How do you remember your experience immediately after the session and later? Did anything change? What happened in between? How do you interpret it now? \\
Context and alternatives & How, if at all, did prior musical experience, the tool, other people or alternative ways of creating matter? \\
Contrasting episode & Can you describe another occasion that differed in a meaningful way? What else might explain the difference? \\
Closure & Is there anything important we have not asked? Would you like to withdraw any part of the interview under the agreed procedure? \\
\bottomrule
\end{tabularx}
\end{table}
Probes will distinguish feelings represented in a song from feelings experienced by a participant. The interviewer will not ask whether a presumed mechanism, such as reappraisal, occurred before eliciting the person's account. A participant may describe an episode without specifying the underlying life event. The closing discussion will check whether the participant wishes to take a break or withdraw any material within the agreed withdrawal procedure.

